% ============================================================
%  Appendix B — Data Sources and Outputs
% ============================================================
\chapter{Data Sources and Outputs}
\label{app:data}

\section{Source Datasets}
\label{app:sources}

All input data were obtained through Google Earth Engine. Table~\ref{tab:gee_full}
gives the full dataset identifiers, native resolutions, and the temporal windows
used.

\begin{table}[H]
\centering
\caption{Google Earth Engine source datasets for the nine input features.}
\label{tab:gee_full}
\small
\begin{tabular}{p{2.3cm} p{5.4cm} p{2.2cm} p{2.6cm}}
\toprule
Feature & GEE dataset & Native res. & Period \\
\midrule
Elevation & \texttt{USGS/SRTMGL1\_003} & 30~m & 2000 \\
Slope, Aspect & \texttt{ee.Terrain} (from SRTM) & 30~m & 2000 \\
TWI & \texttt{WWF/HydroSHEDS/15ACC} + SRTM & 15~arc-sec & static \\
NDVI, NDWI & \texttt{LANDSAT/LC08/C02/T1\_L2} & 30~m & 2015--2023 \\
Rainfall & \texttt{UCSB-CHG/CHIRPS/DAILY} & 0.05$^\circ$ & 2015--2023 \\
Soil texture & \texttt{OpenLandMap/SOL/SOL\_TEXTURE-CLASS\_USDA-TT\_M/v02} & 250~m & static \\
LST & \texttt{MODIS/061/MOD11A1} & 1~km & 2015--2023 \\
\bottomrule
\end{tabular}
\end{table}

All bands were resampled to a common 30~m grid in EPSG:4326 on export, and
reprojected to UTM Zone~42N (EPSG:32642) for area and distance calculations. The
exported raster is $10{,}062 \times 6{,}884$ pixels (about 69.2~million cells), of
which \npixels{} (53.4\%) lie within the five district boundaries and carry valid
data.

\section{Study Area Specification}
\label{app:area}

\begin{table}[H]
\centering
\caption{Study area parameters.}
\label{tab:area_spec}
\begin{tabular}{ll}
\toprule
Parameter & Value \\
\midrule
Districts & Chakwal, Rawalpindi, Attock, Jhelum, Mianwali \\
Approx.\ area & \SI{25000}{\kilo\meter\squared} \\
Latitude range & \ang{32.16}--\ang{34.02}~N \\
Longitude range & \ang{71.11}--\ang{73.82}~E \\
Export CRS & EPSG:4326 \\
Analysis CRS & EPSG:32642 (UTM Zone 42N) \\
Resolution & 30~m \\
Valid pixels & \npixels{} \\
\bottomrule
\end{tabular}
\end{table}

\section{Model Outputs}
\label{app:outputs}

The principal geospatial outputs of the study are listed in
Table~\ref{tab:outputs}. The classified and probability rasters are the
deliverables intended for downstream GIS use.

\begin{table}[H]
\centering
\caption{Principal model outputs.}
\label{tab:outputs}
\small
\begin{tabular}{p{5.2cm} p{8cm}}
\toprule
Output & Description \\
\midrule
\texttt{gw\_prediction\_xgboost.tif} & Classified groundwater potential (0 = Low, 1 = Medium, 2 = High), 30~m. \\
\texttt{gw\_probability\_high.tif}   & Continuous probability of the High class, 30~m. \\
\texttt{gw\_potential\_labels.tif}   & Weighted-overlay reference labels. \\
\texttt{best\_model.joblib}          & Trained XGBoost model. \\
\texttt{metrics.json}                & Full accuracy report and per-class metrics. \\
\bottomrule
\end{tabular}
\end{table}

\section{Predicted Distribution by District}
\label{app:district}

Table~\ref{tab:district_high} reports the share of each district classified as
High potential, which forms the basis of the district-gradient validation in
Section~\ref{sec:validation_results}. The ordering follows the expected
hydrogeological gradient from the wetter, alluvial north to the drier, hard-rock
south (Spearman $\rho = 1.00$).

\begin{table}[H]
\centering
\caption{Predicted High-potential fraction by district.}
\label{tab:district_high}
\begin{tabular}{lc}
\toprule
District & High potential (\%) \\
\midrule
Rawalpindi & 61.3 \\
Attock     & 32.5 \\
Jhelum     & 20.8 \\
Mianwali   & 5.4 \\
Chakwal    & 4.0 \\
\bottomrule
\end{tabular}
\end{table}
